\section{Method}

PaperToSkill represents each converted paper as a skill package. The main file
is \texttt{SKILL.md}; optional references include a source map that links
generated skill bullets to source note sections and line spans. The schema
contains identity, central contribution, use conditions, required artifacts,
workflow steps, validation checks, failure cases, transfer notes, and source
notes that distinguish source-backed material from inferred guidance.

The deterministic extractor reads a source-anchored paper note, normalizes
Markdown sections and list items, collects candidates from abstract, methods,
experiments, limitations, and transfer sections, then emits a compact skill.
Candidate limits keep the artifact within a practical context budget while
preserving method, validation, and failure coverage. The extractor also writes a
\texttt{references/source\_map.json} file so later audits can inspect which
source section supports each generated instruction.

The text-to-note scaffold is an earlier preprocessing step. It operates on
newline-numbered extracted text, detects broad paper regions, selects line
windows for methods, experiments, and limitations, and preserves line anchors.
For two-column extracted text, it keeps raw line spacing and selects the
keyword-bearing column while retaining the original source line range. The
output is explicitly marked as an audit scaffold.

\begin{figure*}[t]
  \centering
  \includegraphics[width=0.96\textwidth]{papertoskill_pipeline.pdf}
  \caption{PaperToSkill follows the flow source-anchored paper $\rightarrow$
  operational note $\rightarrow$ portable skill $\rightarrow$ evaluation.
  Evaluation separates deterministic local checks of package structure,
  grounding, and transfer coverage from the external-evidence layer for
  cross-model real-reuse and independent human assessment.}
  \label{fig:papertoskill-pipeline}
\end{figure*}

\section{Experimental Setup}

We evaluate four real agent-method papers: AI Scientist-v2, Reflexion, AIDE,
and Toolformer. For each paper, we create a curated source-anchored note from
the extracted paper text, generate a PaperToSkill skill, and compare it against
two context baselines: a generic summary and an abstract-only context. For
transfer ablation, we compare the full skill against the same skill with
\texttt{Transfer Notes} removed and against the generic summary.

The metrics are deterministic: skill rubric score, context coverage,
source-span validation, offline transfer-readiness, word count under a
1200-word compactness budget, character-based token/cost proxy, and a local
tokenizer-aware input proxy using \texttt{o200k\_base}. For saved
model-ablation responses, we additionally compute a local output-token proxy.
These metrics are reproducible gates. They do not replace live agent execution,
real billing records, invoice evidence, or success-per-dollar accounting. They
are complemented by a bounded single-reviewer, six-criterion internal
annotation summary with 24/24 scored paper-by-criterion cells, a mean score of
4.75/5, mean confidence 0.946, and four discussion-marked rows. The annotation
supports an internal semantic-fidelity audit, not independent human validation
or live downstream task success.
The current package uses local token accounting over input-token and
saved-response output-token proxies. We also
prepare model-ablation prompt packets for Claude Opus 4.8, a GPT-family slot
requested as GPT 5.5, and a DeepSeek slot, and include a protocol-aware runner
and response evaluator. The current two-case protocol has 6/6 saved and scored
rows across Claude, GPT-family, and DeepSeek slots. This is saved-response
output-contract evidence only; it is not live downstream task success,
provider billing, or a broad model-quality result. Provider/model availability
failures are logged separately from model-quality results.

The primary downstream experiment is the real-reuse protocol in
Table~\ref{tab:real-reuse-main}. It contains eight original-style
\emph{paper-task} rows across AIDE, SWE-agent, Reflexion, and
SnapATAC2~\cite{zhang2024snapatac2}. Each
row fixes the source paper, task input/output shape, metric family, reported
reference boundary, and the two main conditions: Summary and PaperToSkill. At
this revision, all eight rows have one GPT-family run with scored Summary and
PaperToSkill outputs. The AIDE rows use an official Kaggle Spaceship Titanic
training file with a locked local validation split; the SnapATAC2 rows use
prepared official miniature fixtures and are scored local dry-run tasks rather
than full paper reproductions. These rows complete the first single-run pass
over the real-reuse table, but they do not by themselves establish aggregate
advantage over Summary.

We also instantiate a small Full Excerpt sanity scaffold over the
pre-registered AIDE-T1, SWE-T1, and SNAP-T1 subset
(Table~\ref{tab:full-excerpt-sanity}). This table is separate from the main
real-reuse comparison: it records the existing Summary/PaperToSkill scores and
local context-token proxies, and the matched Full Excerpt runs remain
auxiliary sanity evidence rather than a main baseline.

\begin{table*}[t]
\centering
\caption{Main real-reuse experiment table. Filled scores come from local raw rows; future reruns or additional rows may be added. Original-paper scores are reported references unless the same data, input/output, metric, budget, and runtime setting are reproduced locally.}
\label{tab:real-reuse-main}
\resizebox{\textwidth}{!}{%
\begin{tabular}{llllllllll}
\toprule
Task & Paper & Domain & Input & Output & Metric & Reference & Summary & PaperToSkill & Status \\
\midrule
AIDE-T1 & AIDE & ML engineering & Kaggle-style dataset + metric & Runnable solution/submission & validation\_score & Reported AIDE ref. & 0.816 & 0.817 & Scored (GPT-family) \\
AIDE-T2 & AIDE & ML engineering & Weak ML script + feedback & Improved script + trajectory & best\_node\_score & Reported AIDE ref. & 0.000 & 0.826 & Scored (GPT-family) \\
SWE-T1 & SWE-agent & Software engineering & Repo issue + tests & Patch + test log & resolved & Reported SWE-agent ref. & 0.000 & 0.000 & Scored (GPT-family) \\
SWE-T2 & SWE-agent & Software engineering & Failing test + repo & Focused patch + verification & tests\_passed & Reported SWE-agent ref. & 0.000 & 1.000 & Scored (GPT-family) \\
REF-T1 & Reflexion & Reasoning / QA & Multi-hop QA + feedback & Final answer + reflection trace & exact\_match\_or\_f1 & Reported Reflexion ref. & 1.000 & 1.000 & Scored (GPT-family) \\
REF-T2 & Reflexion & Decision / programming & Failed attempt + checker feedback & Corrected second attempt & second\_attempt\_success & Reported Reflexion ref. & 1.000 & 1.000 & Scored (GPT-family) \\
SNAP-T1 & SnapATAC2 & Single-cell omics & Small single-cell dataset & Pipeline + embedding artifacts & runtime\_memory\_quality & Reported SnapATAC2 ref. & 0.000 & 0.500 & Scored (GPT-family) \\
SNAP-T2 & SnapATAC2 & Single-cell omics & Single-cell labels/proxy task & Clustering/marker artifacts & ari\_nmi\_runtime\_memory & Reported SnapATAC2 ref. & 0.200 & 0.400 & Scored (GPT-family) \\
\bottomrule
\end{tabular}%
}
\end{table*}

\begin{table*}[t]
\centering
\caption{Real-reuse failure-boundary analysis derived from the same scored raw rows as Table~\ref{tab:real-reuse-main}. This table explains first-pass boundary modes; it does not add new task-success evidence.}
\label{tab:real-reuse-failure-analysis}
\resizebox{\textwidth}{!}{%
\begin{tabular}{lllll}
\toprule
Task & Summary Outcome & PaperToSkill Outcome & Boundary Mode & Contract Implication \\
\midrule
AIDE-T1 & 0.816; success & 0.817; success & Solved by both & harder task slice or stricter scorer \\
AIDE-T2 & 0.000; timeout & 0.826; success & PaperToSkill-only success & preserve patch/tool-use constraints \\
SWE-T1 & 0.000; patch apply failed & 0.000; patch apply failed & Patch application & patch-format and apply-check contract \\
SWE-T2 & 0.000; patch apply failed & 1.000; success & PaperToSkill-only success & preserve patch/tool-use constraints \\
REF-T1 & 1.000; success & 1.000; success & Solved by both & harder task slice or stricter scorer \\
REF-T2 & 1.000; success & 1.000; success & Solved by both & harder task slice or stricter scorer \\
SNAP-T1 & 0.000; artifact parse failed & 0.500; artifact incomplete & Artifact completion & required artifact and metric manifest \\
SNAP-T2 & 0.200; artifact incomplete & 0.400; artifact incomplete & Artifact completion & required artifact and metric manifest \\
\bottomrule
\end{tabular}%
}
\end{table*}

\begin{table*}[t]
\centering
\caption{SWE-T1 shared-source-context follow-up. The first-pass SWE-T1 row remains the main-table evidence; this paired follow-up exposes the same locked SQLFluff source slice to both conditions and is diagnostic rather than a replacement.}
\label{tab:swe-t1-source-context-followup}
\resizebox{\textwidth}{!}{%
\begin{tabular}{llllllllllll}
\toprule
Task & Condition & First Run & First Score & First Failure & First Patch & Follow-up Run & Follow-up Score & Follow-up Patch & Follow-up Test & Follow-up Failure & Interpretation \\
\midrule
SWE-T1 & Summary & phase97\_gpt\_swe\_t1\_summary\_retry & 0.000 & patch\_apply\_failed & No & phase107\_gpt\_swe\_t1\_source\_context\_followup & 0.000 & Yes & No & test\_command\_failed & Source context fixed patch application; hidden test still failed \\
SWE-T1 & PaperToSkill & phase97\_gpt\_swe\_t1\_real\_reuse & 0.000 & patch\_apply\_failed & No & phase107\_gpt\_swe\_t1\_source\_context\_followup & 0.000 & Yes & No & test\_command\_failed & Source context fixed patch application; hidden test still failed \\
\bottomrule
\end{tabular}%
}
\end{table*}

\begin{table*}[t]
\centering
\caption{SWE-T1 issue-aligned follow-up. The first-pass SWE-T1 row remains the main-table evidence; this diagnostic uses the pre-registered issue-aligned scorer and shows that both Summary and PaperToSkill pass under the revised contract, so it is not evidence of PaperToSkill advantage.}
\label{tab:swe-t1-issue-aligned-followup}
\resizebox{\textwidth}{!}{%
\begin{tabular}{llllllllllll}
\toprule
Task & Condition & Main Run & Main Score & Phase107 Run & Phase107 Score & Phase107 Test & Issue Run & Issue Score & Issue Test & Issue Test Patch & Interpretation \\
\midrule
SWE-T1 & Summary & phase97\_gpt\_swe\_t1\_summary\_retry & 0.000 & phase107\_gpt\_swe\_t1\_source\_context\_followup & 0.000 & No & phase110\_gpt\_swe\_t1\_issue\_aligned\_followup & 1.000 & Yes & No & Issue-aligned scorer passes both conditions; no PaperToSkill advantage \\
SWE-T1 & PaperToSkill & phase97\_gpt\_swe\_t1\_real\_reuse & 0.000 & phase107\_gpt\_swe\_t1\_source\_context\_followup & 0.000 & No & phase110\_gpt\_swe\_t1\_issue\_aligned\_followup & 1.000 & Yes & No & Issue-aligned scorer passes both conditions; no PaperToSkill advantage \\
\bottomrule
\end{tabular}%
}
\end{table*}

\begin{table*}[t]
\centering
\caption{SnapATAC2 executable-artifact follow-up. The main SNAP rows in Table~\ref{tab:real-reuse-main} remain unchanged; this paired diagnostic executes a pre-registered controlled scaffold over the same miniature fixtures to test whether concrete artifacts plus runtime/memory records satisfy the existing scorer.}
\label{tab:snapatac2-executable-followup}
\resizebox{\textwidth}{!}{%
\begin{tabular}{lllllll}
\toprule
Task & Condition & Run & Score & Success & Runtime Seconds & Peak Memory MB \\
\midrule
SNAP-T1 & summary & phase108\_snapatac2\_executable\_artifact\_followup & 1.0 & True & 1.296368 & 26.378289 \\
SNAP-T1 & papertoskill & phase108\_snapatac2\_executable\_artifact\_followup & 1.0 & True & 1.209319 & 26.338634 \\
SNAP-T2 & summary & phase108\_snapatac2\_executable\_artifact\_followup & 1.0 & True & 1.75799 & 24.648565 \\
SNAP-T2 & papertoskill & phase108\_snapatac2\_executable\_artifact\_followup & 1.0 & True & 1.756513 & 24.758577 \\
\bottomrule
\end{tabular}%
}
\end{table*}

\begin{table*}[t]
\centering
\caption{SnapATAC2 SNAP-T1 executable-candidate follow-up. The main SNAP rows in Table~\ref{tab:real-reuse-main} remain unchanged; this diagnostic executes model-generated phase112 candidate scripts under the pre-registered executable-candidate contract and shows contract closure for both conditions, not PaperToSkill advantage.}
\label{tab:snapatac2-executable-candidate-followup}
\resizebox{\textwidth}{!}{%
\begin{tabular}{lllllll}
\toprule
Task & Condition & Run & Score & Success & Runtime Seconds & Peak Memory MB \\
\midrule
SNAP-T1 & summary & phase112\_gpt\_snapatac2\_executable\_candidate\_t1 & 1.0 & True & 9.362002 & 0.0 \\
SNAP-T1 & papertoskill & phase112\_gpt\_snapatac2\_executable\_candidate\_t1 & 1.0 & True & 2.010309 & 164.615443 \\
\bottomrule
\end{tabular}%
}
\end{table*}

\begin{table*}[t]
\centering
\caption{Auxiliary Full Excerpt sanity check over the pre-registered subset. Scores come from matched real-reuse rows where available; token columns are local whitespace context proxies, not provider billing or output-token costs.}
\label{tab:full-excerpt-sanity}
\resizebox{\textwidth}{!}{%
\begin{tabular}{llllllllll}
\toprule
Task & Paper & Metric & Summary & PaperToSkill & Full Excerpt & Summary Tokens & PaperToSkill Tokens & Full Excerpt Tokens & Status \\
\midrule
AIDE-T1 & AIDE & validation\_score & 0.816 & 0.817 & 0.000 & 121 & 878 & 7366 & Scored (GPT-family) \\
SWE-T1 & SWE-agent & resolved & 0.000 & 0.000 & 0.000 & 89 & 1173 & 42048 & Scored (GPT-family) \\
SNAP-T1 & SnapATAC2 & runtime\_memory\_quality & 0.000 & 0.500 & 0.250 & 72 & 1069 & 10297 & Scored (GPT-family) \\
\bottomrule
\end{tabular}%
}
\end{table*}

\begin{table*}[t]
\centering
\caption{Main deterministic/offline PaperToSkill results. Coverage scores are task-rubric scores; transfer readiness is an offline artifact-readiness metric.}
\label{tab:main-results}
\begin{tabular}{lrrrrrrrr}
\toprule
Paper & Rubric & Skill & Summary & Abstract & $\Delta$Sum. & Transfer & Support & Words \\
\midrule
AI Scientist-v2 & 20/20 & 7.867/9 & 1.733/9 & 1.2/9 & 6.134 & 10/10 & 0.938 & 782 \\
Reflexion & 20/20 & 8.267/9 & 3.483/9 & 2.533/9 & 4.784 & 10/10 & 1.000 & 479 \\
AIDE & 20/20 & 9.1/10 & 1.916/10 & 1.333/10 & 7.184 & 10/10 & 1.000 & 927 \\
Toolformer & 20/20 & 8.9/10 & 2.5/10 & 1.534/10 & 6.400 & 10/10 & 1.000 & 943 \\
\bottomrule
\end{tabular}
\end{table*}

\begin{table}[t]
\centering
\caption{Transfer-note ablation. Removing only \texttt{Transfer Notes} lowers offline readiness across all curated skills.}
\label{tab:transfer-ablation}
\begin{tabular}{lrr}
\toprule
Paper & Full & No Transfer \\
\midrule
AI Scientist-v2 & 10/10 & 7.6/10 \\
Reflexion & 10/10 & 7.6/10 \\
AIDE & 10/10 & 7.6/10 \\
Toolformer & 10/10 & 7.6/10 \\
\bottomrule
\end{tabular}
\end{table}

\begin{table*}[t]
\centering
\caption{Local tokenizer-aware context-size proxy using the \texttt{o200k\_base} tokenizer. Counts are not provider bills, output-token costs, or success-per-dollar measurements.}
\label{tab:cost-proxy}
\begin{tabular}{lrrrr}
\toprule
Paper & Full Paper Tokens & Skill Tokens & Skill/Full & Reduction \\
\midrule
AI Scientist-v2 & 45212 & 1079 & 0.0239 & 97.61\% \\
Reflexion & 16414 & 703 & 0.0428 & 95.72\% \\
AIDE & 13312 & 1285 & 0.0965 & 90.35\% \\
Toolformer & 20365 & 1255 & 0.0616 & 93.84\% \\
\bottomrule
\end{tabular}
\end{table*}

\begin{table*}[t]
\centering
\caption{Automatic extracted-text note scaffolds are evaluated separately from curated-note cases.}
\label{tab:auto-note}
\begin{tabular}{llrrrrr}
\toprule
Paper & Input & Rubric & Coverage & Transfer & Support & Words \\
\midrule
Toolformer & Curated note & 20/20 & 8.9/10 & 10/10 & 1.000 & 943 \\
Toolformer & Auto note scaffold & 20/20 & 9.3/10 & 10/10 & 1.000 & 1179 \\
AIDE & Curated note & 20/20 & 9.1/10 & 10/10 & 1.000 & 927 \\
AIDE & Auto note scaffold & 20/20 & 8.467/10 & 9.5/10 & 1.000 & 998 \\
\bottomrule
\end{tabular}
\end{table*}


\begin{figure*}[t]
  \centering
  \includegraphics[width=0.96\textwidth]{external_evidence_design.pdf}
  \caption{External evidence design with two tracks: cross-model live
  real-reuse evaluation and expanded independent human evaluation.}
  \label{fig:external-evidence-design}
\end{figure*}

\newcommand{\pendingevidencevalue}{\texttt{-{}-}}

\begin{table*}[t]
\centering
\caption{Cross-model live real-reuse evaluation.}
\label{tab:external-live-real-reuse}
\resizebox{\textwidth}{!}{%
\begin{tabular}{lllllll}
\toprule
Model Family & Task & Seeds & Summary & PaperToSkill & Delta & Status \\
\midrule
GPT-family & AIDE-T2 & \pendingevidencevalue & \pendingevidencevalue & \pendingevidencevalue & \pendingevidencevalue & Pending external evidence \\
GPT-family & SWE-T2 & \pendingevidencevalue & \pendingevidencevalue & \pendingevidencevalue & \pendingevidencevalue & Pending external evidence \\
GPT-family & REF-T2 & \pendingevidencevalue & \pendingevidencevalue & \pendingevidencevalue & \pendingevidencevalue & Pending external evidence \\
Claude-family & AIDE-T2 & \pendingevidencevalue & \pendingevidencevalue & \pendingevidencevalue & \pendingevidencevalue & Pending external evidence \\
Claude-family & SWE-T2 & \pendingevidencevalue & \pendingevidencevalue & \pendingevidencevalue & \pendingevidencevalue & Pending external evidence \\
Claude-family & REF-T2 & \pendingevidencevalue & \pendingevidencevalue & \pendingevidencevalue & \pendingevidencevalue & Pending external evidence \\
DeepSeek-family & AIDE-T2 & \pendingevidencevalue & \pendingevidencevalue & \pendingevidencevalue & \pendingevidencevalue & Pending external evidence \\
DeepSeek-family & SWE-T2 & \pendingevidencevalue & \pendingevidencevalue & \pendingevidencevalue & \pendingevidencevalue & Pending external evidence \\
DeepSeek-family & REF-T2 & \pendingevidencevalue & \pendingevidencevalue & \pendingevidencevalue & \pendingevidencevalue & Pending external evidence \\
\bottomrule
\end{tabular}%
}
\end{table*}

\begin{table*}[t]
\centering
\caption{Expanded independent human evaluation.}
\label{tab:external-independent-human}
\resizebox{\textwidth}{!}{%
\begin{tabular}{lllllll}
\toprule
Criterion & Artifacts & Raters & Mean & Agreement & 95\% CI & Status \\
\midrule
Central contribution fidelity & \pendingevidencevalue & \pendingevidencevalue & \pendingevidencevalue & \pendingevidencevalue & \pendingevidencevalue & Pending external evidence \\
Operational workflow fidelity & \pendingevidencevalue & \pendingevidencevalue & \pendingevidencevalue & \pendingevidencevalue & \pendingevidencevalue & Pending external evidence \\
Validation and evidence fidelity & \pendingevidencevalue & \pendingevidencevalue & \pendingevidencevalue & \pendingevidencevalue & \pendingevidencevalue & Pending external evidence \\
Failure and limitation fidelity & \pendingevidencevalue & \pendingevidencevalue & \pendingevidencevalue & \pendingevidencevalue & \pendingevidencevalue & Pending external evidence \\
Source grounding & \pendingevidencevalue & \pendingevidencevalue & \pendingevidencevalue & \pendingevidencevalue & \pendingevidencevalue & Pending external evidence \\
Transfer boundary discipline & \pendingevidencevalue & \pendingevidencevalue & \pendingevidencevalue & \pendingevidencevalue & \pendingevidencevalue & Pending external evidence \\
\bottomrule
\end{tabular}%
}
\end{table*}

